% AA-04 | Algebra abstrakcyjna | Iteracja 02
\chapter{Twierdzenie Cauchy'ego i trzy twierdzenia Sylowa}
\label{aa:ch:sylow}

\section{Od twierdzenia Lagrange'a do problemu istnienia podgrup}
Twierdzenie Lagrange'a mówi, że jeśli $H\le G$ i $G$ jest
skończona, to $|H|$ dzieli $|G|$. Nie wynika z niego, że
każdy dzielnik $|G|$ jest rzędem podgrupy: grupa $A_4$ rzędu
$12$ nie ma podgrupy rzędu $6$.
Nasuwa się jednak słabsze pytanie: jeśli liczba pierwsza $p$
dzieli $|G|$, czy istnieje w $G$ element rzędu $p$?
Odpowiedź jest twierdząca. Co więcej, jeśli $p^a$ jest
największą potęgą $p$ dzielącą $|G|$, to istnieje nawet
podgrupa rzędu $p^a$. Dowiedziemy obu faktów bez zakładania
jakiejkolwiek klasyfikacji grup skończonych.

\begin{definicja}[Podgrupa Sylowa]
\label{aa:def:sylow}
Niech $G$ będzie skończoną grupą, $p$ liczbą pierwszą,
a $|G|=p^a m$, gdzie $a$ jest nieujemną liczbą całkowitą,
$m$ jest dodatnią liczbą całkowitą i $p\nmid m$.
\emph{Podgrupą $p$-Sylowa} grupy $G$ nazywamy podgrupę
rzędu $p^a$. Zbiór wszystkich takich podgrup oznaczamy
$\operatorname{Syl}_p(G)$, a ich liczbę przez $n_p(G)$.
Jeżeli $p\nmid|G|$, jedyną podgrupą $p$-Sylowa jest $\{e\}$.
\end{definicja}

\begin{uwaga}[Dlaczego rząd $p^a$ jest maksymalny?]
Jeżeli podgrupa $Q\le G$ ma rząd będący potęgą $p$,
to twierdzenie Lagrange'a daje $|Q|\mid p^a m$.
Stąd $|Q|=p^b$ dla pewnego $b\le a$. Podgrupa Sylowa
osiąga największy możliwy rząd wśród takich podgrup.
Samo pojęcie jest sensowne przed dowodem istnienia,
lecz dowód istnienia jest główną treścią pierwszego twierdzenia Sylowa.
\end{uwaga}

\section{Twierdzenie Cauchy'ego}
\begin{lemat}[Orbity działania grupy rzędu $p$]
\label{aa:lem:orbit-cp}
Niech $p$ będzie liczbą pierwszą, a cykliczna grupa $C_p$
rzędu $p$ niech działa na skończonym zbiorze $X$.
Każda orbita ma dokładnie $1$ albo $p$ elementów.
Jeżeli $F$ oznacza zbiór punktów ustalanych przez wszystkie
elementy $C_p$, to $|X|\equiv |F|\pmod p$.
\end{lemat}
\begin{proof}
Stabilizator każdego punktu jest podgrupą $C_p$.
Z twierdzenia Lagrange'a jego rząd dzieli $p$,
więc wynosi $1$ albo $p$.
Z twierdzenia o orbicie i stabilizatorze
rozmiar orbity to odpowiednio $p$ albo $1$.
Orbita jest jednoelementowa wtedy i tylko wtedy,
gdy wszystkie elementy grupy ustalają jej jedyny punkt.
Zbiór $X$ rozpada się zatem na $|F|$ orbit jednoelementowych
i pewną liczbę orbit $p$-elementowych.
Licząc elementy, dostajemy $|X|=|F|+kp$
przy pewnym $k\ge0$, co jest równoważne kongruencji.
\end{proof}

\begin{twierdzenie}[Cauchy'ego]
\label{aa:tw:cauchy}
Niech $G$ będzie grupą skończoną, a $p$ liczbą pierwszą
dzielącą $|G|$. Wtedy w $G$ istnieje element rzędu $p$.
W konsekwencji $G$ zawiera podgrupę rzędu $p$.
\end{twierdzenie}
\begin{proof}
Rozważmy zbiór
\[
 X=\{(g_1,\ldots,g_p)\in G^p:g_1g_2\cdots g_p=e\}.
\]
Pierwsze $p-1$ współrzędnych możemy wybrać zupełnie dowolnie.
Ostatnia współrzędna jest wówczas jednoznacznie wyznaczona:
$g_p=(g_1\cdots g_{p-1})^{-1}$.
Wobec tego $|X|=|G|^{p-1}$, a ponieważ $p\mid |G|$,
liczba $|X|$ jest podzielna przez $p$.

Niech generator cyklicznej grupy $C_p$ działa na $X$
przez cykliczne przesunięcie współrzędnych:
$T(g_1,\ldots,g_p)=(g_2,\ldots,g_p,g_1)$.
Musimy sprawdzić, że $T$ zachowuje $X$, zwłaszcza że $G$
nie musi być przemienna. Jeśli $g_1\cdots g_p=e$, to
\[
 g_2\cdots g_p g_1
 =g_1^{-1}(g_1\cdots g_p)g_1
 =g_1^{-1}eg_1=e.
\]
Zatem $T:X\to X$ jest dobrze określone.
Po $p$ przesunięciach wracamy do początkowej krotki,
więc $T^p=\operatorname{id}_X$ i potęgi $T$ wyznaczają
działanie $C_p$ na $X$ (nawet gdy to działanie nie jest wierne).

Punkty ustalone przez generator $T$ to dokładnie krotki
$(g,g,\ldots,g)$, dla których $g^p=e$.
Rzeczywiście równość $T(g_1,\ldots,g_p)=(g_1,\ldots,g_p)$
wymusza $g_1=\cdots=g_p=g$,
a warunek przynależności do $X$ daje $g^p=e$.
Ponieważ $C_p$ jest generowana przez $T$,
są to również punkty stałe całego działania.
Oznaczmy ich liczbę przez $r$.
Z lematu o orbitach $r\equiv|X|\equiv0\pmod p$.

Jedną z tych krotek jest $(e,\ldots,e)$, zatem $r\ge1$.
Jako dodatnia wielokrotność $p$ liczba $r$ musi być
co najmniej $p$, więc jest też krotka $(g,\ldots,g)$
z $g\ne e$ i $g^p=e$.
Rząd $g$ dzieli $p$ na mocy własności rzędów elementów,
a ponieważ $g\ne e$, nie jest równy $1$.
Zatem $\operatorname{ord}(g)=p$.
Podgrupa $\langle g\rangle$ składa się z dokładnie $p$ elementów.
\end{proof}

\begin{przyklad}[Jak twierdzenie Cauchy'ego działa w $S_3$?]
Grupa $S_3$ ma rząd $6=2\cdot3$.
Dla $p=2$ twierdzenie daje element rzędu $2$,
na przykład transpozycję $(1\,2)$.
Dla $p=3$ daje element rzędu $3$,
na przykład $(1\,2\,3)$.
Twierdzenie nie mówi, że istnieje element rzędu $6$:
w $S_3$ żaden element nie ma tego rzędu, mimo że
istnieją elementy rzędów $2$ i $3$.
Przyczyna jest strukturalna: iloczyn dwóch takich elementów
nie musi mieć rzędu będącego iloczynem ich rzędów,
gdyż elementy na ogół nie komutują.
\end{przyklad}

\begin{przyklad}[Liczba punktów stałych w konstrukcji Cauchy'ego]
\label{aa:prz:cauchy-z6}
Niech $G=(\mathbb Z/6\mathbb Z,+)$.
W dowodzie twierdzenia Cauchy'ego warunek iloczynu równego
$e$ staje się warunkiem sumy współrzędnych równej $[0]$.
Dla $p=2$ mamy $|X|=6$, a punkty stałe przesunięcia
$(x_1,x_2)\mapsto(x_2,x_1)$ są postaci $([a],[a])$
z $2[a]=[0]$. Są dokładnie dwa: $a=0$ oraz $a=3$.
Pozostałe $6-2=4$ krotki tworzą dwie orbity dwuelementowe,
więc istotnie $|X|=2+2\cdot2$.
Dla $p=3$ mamy $|X|=6^2=36$.
Punkty stałe cyklicznego przesunięcia są postaci
$([a],[a],[a])$ z $3[a]=[0]$, czyli $a\in\{0,2,4\}$.
Istnieją zatem trzy punkty stałe i
$(36-3)/3=11$ orbit trójelementowych.
W obu przypadkach dodatkowe nietrywialne rozwiązanie
równania $p[a]=[0]$ daje element rzędu $p$.
Dla ogólnej grupy nieprzemiennej rachunek przebiega
analogicznie, lecz należy użyć iloczynów zamiast sum.
\end{przyklad}

\begin{wniosek}[Centralny element rzędu $p$ w $p$-grupie]
\label{aa:wn:central-p}
Niech $G$ będzie skończoną grupą rzędu $p^a$, gdzie
$p$ jest liczbą pierwszą i $a\ge1$.
Wtedy w centrum $Z(G)$ istnieje element rzędu $p$.
W szczególności $G$ ma normalną podgrupę rzędu $p$.
\end{wniosek}
\begin{proof}
Z Twierdzenia~\ref{aa:tw:centrum-pgrupy} wiemy,
że $p\mid|Z(G)|$. Centrum jest skończoną grupą,
więc twierdzenie Cauchy'ego zastosowane do $Z(G)$
daje element $z$ rzędu $p$.
Podgrupa $\langle z\rangle$ ma rząd $p$.
Każda jej potęga należy do centrum, więc dla wszystkich
$g\in G$ mamy $g\langle z\rangle g^{-1}=\langle z\rangle$:
podgrupa jest normalna.
\end{proof}

\section{Lemat kombinatoryczny przygotowujący dowód Sylowa}
Pierwsze twierdzenie Sylowa wyprowadzimy z działania
na zbiorze podzbiorów grupy określonej liczności.
Kluczowe będzie wykazanie, że pewien współczynnik
dwumianowy nie dzieli się przez $p$. Poniższy lemat
uzasadnia to bez odwołań do nieudowodnionych wyników.

\begin{lemat}[Współczynnik dwumianowy modulo $p$]
\label{aa:lem:dwumian-p}
Niech $p$ będzie liczbą pierwszą, $a\ge0$ i $m\ge1$,
gdzie $p\nmid m$. Dla $n=p^a m$ zachodzi
\[
 \binom{n}{p^a}\equiv m\pmod p.
\]
W szczególności $p$ nie dzieli $\binom{n}{p^a}$.
\end{lemat}
\begin{proof}
Pracujemy w pierścieniu wielomianów
$(\mathbb Z/p\mathbb Z)[T]$; współczynniki traktujemy
jako klasy reszt modulo $p$.
Dla $1\le j\le p-1$ zachodzi $p\mid\binom pj$.
Aby to wykazać, pomnóżmy wzór
$j\binom pj=p\binom{p-1}{j-1}$.
Prawa strona jest podzielna przez $p$.
Ponieważ $p$ jest pierwsza i $p\nmid j$,
to $p\mid\binom pj$ (wynika to również z tożsamości
Bézouta, gdyż istnieją $u,v$ z $uj+vp=1$).
Po redukcji współczynników dwumian Newtona daje zatem
$(1+T)^p=1+T^p$ w naszym pierścieniu wielomianów.

Zastosujemy tę tożsamość wielokrotnie.
Dla $a=0$ mamy $(1+T)^{p^0}=1+T$.
Jeśli $(1+T)^{p^a}=1+T^{p^a}$, to podnosząc obie
strony do potęgi $p$ otrzymujemy
$(1+T)^{p^{a+1}}=(1+T^{p^a})^p=1+T^{p^{a+1}}$.
Indukcja dowodzi wzoru dla wszystkich $a\ge0$.

Wobec tego
$(1+T)^n=((1+T)^{p^a})^m=(1+T^{p^a})^m$
w pierścieniu $(\mathbb Z/p\mathbb Z)[T]$.
Porównajmy współczynniki przy $T^{p^a}$.
Po lewej jest to klasa $\binom{n}{p^a}$.
Po prawej wzór dwumianowy
$\sum_{j=0}^m\binom mj T^{jp^a}$ pokazuje,
że jedynym składnikiem o żądanym stopniu jest ten dla $j=1$,
zatem współczynnik wynosi $\binom m1=m$.
Otrzymujemy kongruencję z tezy. Ponieważ $p\nmid m$,
współczynnik dwumianowy nie dzieli się przez $p$.
\end{proof}

\section{Pierwsze twierdzenie Sylowa: istnienie}
\begin{twierdzenie}[Pierwsze twierdzenie Sylowa]
\label{aa:tw:sylow1}
Niech $G$ będzie grupą skończoną, $p$ liczbą pierwszą,
a $|G|=p^a m$, gdzie $a$ jest nieujemną liczbą całkowitą,
$m$ jest dodatnią liczbą całkowitą i $p\nmid m$.
Wtedy istnieje podgrupa $P\le G$ o rzędzie $p^a$,
czyli podgrupa $p$-Sylowa.
\end{twierdzenie}
\begin{proof}
Dla $a=0$ wystarczy wziąć $P=\{e\}$.
Załóżmy więc $a\ge1$ i oznaczmy $r=p^a$.
Niech $\mathcal X$ będzie zbiorem wszystkich $r$-elementowych
podzbiorów grupy $G$. Ponieważ $r\le|G|$, jest to zbiór niepusty,
a jego liczebność wynosi
$|\mathcal X|=\binom{|G|}{r}=\binom{p^a m}{p^a}$.
Na mocy poprzedniego lematu $p\nmid|\mathcal X|$.

Grupa $G$ działa na $\mathcal X$ przez przesunięcia z lewej:
$g\cdot A=\{ga:a\in A\}$.
To działanie jest dobrze określone, bo lewa translacja
$G\to G$, $a\mapsto ga$, jest bijekcją, więc zachowuje
liczbę elementów podzbioru.
Neutralny element ustala każdy podzbiór,
a $g\cdot(h\cdot A)=(gh)\cdot A$.
Zatem faktycznie mamy działanie grupy.

Orbity tworzą rozkład $\mathcal X$.
Gdyby każda orbita miała liczebność podzielną przez $p$,
cała suma ich liczebności, czyli $|\mathcal X|$,
też byłaby podzielna przez $p$. To przeczy lematowi.
Istnieje więc orbita $\mathcal O=G\cdot A$
o liczebności niepodzielnej przez $p$.

Niech $H=G_A=\{g\in G:gA=A\}$ będzie stabilizatorem
wybranego podzbioru $A$.
Z twierdzenia o orbicie i stabilizatorze
$|\mathcal O|=[G:H]=|G|/|H|$.
Liczba $|\mathcal O|$ nie dzieli się przez $p$, a
$|G|=p^a m$. Stąd $p^a$ musi dzielić $|H|$:
wszystkie $a$ czynniki $p$ występujące w $|G|$
pozostają w $|H|$, skoro w ilorazie $|G|/|H|$
nie występuje żaden czynnik $p$.

Z drugiej strony $H$ działa na samym zbiorze $A$ przez
mnożenie z lewej: dla $h\in H$ mamy $hA=A$.
Działanie to jest \emph{swobodne}: jeżeli $h\cdot x=x$
dla pewnego $x\in A$, to $hx=x$ w grupie $G$,
a po pomnożeniu przez $x^{-1}$ otrzymujemy $h=e$.
Zatem stabilizator każdego punktu $x\in A$ jest trywialny.
Na mocy twierdzenia o orbicie i stabilizatorze
każda orbita tego działania ma dokładnie $|H|$ elementów.
Ponieważ orbity rozkładają $A$ na rozłączne części,
$|H|$ dzieli $|A|=p^a$.

Połączywszy oba wnioski, mamy $p^a\mid|H|$ oraz
$|H|\mid p^a$; dlatego $|H|=p^a$.
Stabilizator $H$ jest więc szukaną podgrupą $p$-Sylowa.
\end{proof}

\begin{przyklad}[Istnienie podgrup Sylowa w $S_3$]
Rząd $S_3$ wynosi $6=2\cdot3$.
Pierwsze twierdzenie Sylowa zapewnia istnienie podgrupy rzędu $2$
i podgrupy rzędu $3$.
Przykładami są
$\langle(1\,2)\rangle=\{e,(1\,2)\}$
oraz $\langle(1\,2\,3)\rangle=\{e,(1\,2\,3),(1\,3\,2)\}$.
W tym przypadku obie podgrupy są cykliczne,
ale ogólnie podgrupy Sylowa nie muszą być cykliczne:
w $A_4$ podgrupa rzędu $4$ może być izomorficzna
z $(\mathbb Z/2\mathbb Z)^2$.
\end{przyklad}

\begin{przyklad}[Wydobycie podgrupy Sylowa ze stabilizatora]
\label{aa:prz:sylow-stabilizator-s3}
Rozpatrzmy dowód istnienia podgrup Sylowa w $G=S_3$.
Dla $p=2$ bierzemy zbiór $\mathcal X$ dwuelementowych
podzbiorów $G$. Ma on $\binom62=15$ elementów,
czyli liczebność nieparzystą.
Wybierzmy $A=\{e,(1\,2)\}$.
Stabilizatorem $A$ w działaniu przez lewe translacje jest
$G_A=A$: jeśli $gA=A$, to $g=ge\in A$; z drugiej strony
każdy element dwuelementowej podgrupy $A$ permutuje $A$
przez mnożenie. Orbita $A$ ma zatem $6/2=3$ elementy.
Są nimi $A$, $(1\,3)A$ i $(2\,3)A$, które są parami
różne, gdyż reprezentują różne lewe warstwy $G/A$.
Znaleźliśmy orbitę rozmiaru niepodzielnego przez $2$,
a jej stabilizator jest poszukiwaną podgrupą $2$-Sylowa.

Dla $p=3$ wybieramy podzbiór
$B=\{e,(1\,2\,3),(1\,3\,2)\}=A_3$.
Zbiór trójelementowych podzbiorów ma $\binom63=20$
elementów, a orbita $B$ składa się z dwóch warstw grupy
$S_3$ względem $A_3$. Jej rozmiar wynosi $2$,
nie dzieli się przez $3$ i stabilizator $G_B=B$
ma rząd $3$. Ważne: nie każda orbita musi mieć liczebność
niepodzielną przez $p$; dowód istnienia zapewnia
co najmniej jedną taką orbitę, a konkretny wybór
$A$ lub $B$ pokazuje mechanizm konstrukcji.
\end{przyklad}

\section{Drugie twierdzenie Sylowa: zawieranie i sprzężoność}
\begin{lemat}[Orbity $p$-grupy na skończonym zbiorze]
\label{aa:lem:orbity-pgrupy}
Niech $Q$ będzie skończoną grupą rzędu $p^b$, gdzie $p$
jest liczbą pierwszą, a $b\ge0$, i niech $Q$ działa
na skończonym zbiorze $X$.
Jeżeli $p\nmid|X|$, to działanie ma punkt stały.
\end{lemat}
\begin{proof}
Orbita punktu $x$ ma rozmiar $[Q:Q_x]$.
Na mocy twierdzenia Lagrange'a jest to potęga $p$,
czyli $1,p,p^2,\ldots$.
Jeśli nie byłoby orbit jednoelementowych, każda orbita
miałaby liczebność podzielną przez $p$,
a więc $|X|$, jako suma ich liczebności,
też dzieliłoby się przez $p$. Otrzymalibyśmy sprzeczność.
Istnieje zatem orbita jednoelementowa,
a jej punkt jest ustalany przez całe $Q$.
\end{proof}

\begin{twierdzenie}[Drugie twierdzenie Sylowa]
\label{aa:tw:sylow2}
Niech $G$ będzie grupą skończoną, $p$ liczbą pierwszą,
a $P\in\operatorname{Syl}_p(G)$.
Dla każdej podgrupy $Q\le G$ o rzędzie będącym potęgą $p$
istnieje $g\in G$ takie, że $g^{-1}Qg\subseteq P$.
W szczególności każda taka podgrupa zawiera się w pewnej
podgrupie $p$-Sylowa, a wszystkie podgrupy $p$-Sylowa
są wzajemnie sprzężone.
\end{twierdzenie}
\begin{proof}
Zapiszmy $|G|=p^a m$ z $p\nmid m$.
Ponieważ $|P|=p^a$, zbiór warstw lewostronnych
$G/P$ ma dokładnie $m$ elementów.
Grupa $Q$ działa na $G/P$ przez mnożenie z lewej:
$q\cdot(gP)=(qg)P$. Poprawność tego działania
wykazaliśmy w Przykładzie~\ref{aa:prz:akcja-kosety}.
Ponieważ $Q$ jest $p$-grupą, a liczba $m$ nie jest
podzielna przez $p$, poprzedni lemat zapewnia warstwę
$gP$ ustalaną przez wszystkie elementy $q\in Q$.
Dla każdego $q\in Q$ mamy więc
$qgP=gP$, czyli $g^{-1}qg\in P$.
Oznacza to dokładnie zawieranie $g^{-1}Qg\subseteq P$.
Po sprzężeniu obustronnie przez $g$ otrzymujemy
$Q\subseteq gPg^{-1}$, a grupa $gPg^{-1}$ jest
podgrupą $p$-Sylowa, gdyż sprzężenie jest izomorfizmem.

Jeśli $Q$ także jest podgrupą $p$-Sylowa, to
$|Q|=|P|=p^a$.
Sprzężona grupa $g^{-1}Qg$ ma tyle samo elementów,
co $Q$. Zawiera się w $P$ i ma taką samą liczebność,
zatem $g^{-1}Qg=P$, czyli $Q=gPg^{-1}$.
Dowodzi to sprzężoności wszystkich podgrup Sylowa.
\end{proof}

\begin{uwaga}[Dlaczego używamy lewych warstw?]
Dla dowolnej podgrupy $P\le G$ przestrzeń $G/P$
jest zbiorem warstw lewostronnych, ale nie musi mieć
struktury grupy. W dowodzie wystarcza działanie $Q$
na \emph{zbiorze} warstw, nie zakładamy normalności $P$.
To celowo słabsza konstrukcja niż grupa ilorazowa.
\end{uwaga}

\section{Trzecie twierdzenie Sylowa: liczba podgrup}
\begin{twierdzenie}[Trzecie twierdzenie Sylowa]
\label{aa:tw:sylow3}
Niech $G$ będzie grupą skończoną i niech
$|G|=p^a m$, gdzie $p$ jest liczbą pierwszą,
$a$ jest nieujemną liczbą całkowitą, a $m$ jest dodatnią
liczbą całkowitą niepodzielną przez $p$.
Wtedy liczba $n_p(G)$ podgrup $p$-Sylowa spełnia
\[
 n_p(G)\mid m,\qquad n_p(G)\equiv1\pmod p.
\]
Ponadto dla dowolnej podgrupy $P\in\operatorname{Syl}_p(G)$
mamy $n_p(G)=[G:N_G(P)]$.
\end{twierdzenie}
\begin{proof}
Pierwsze twierdzenie Sylowa zapewnia, że
$\operatorname{Syl}_p(G)$ jest niepusty.
Grupa $G$ działa na zbiorze podgrup Sylowa przez sprzężenie:
$g\cdot Q=gQg^{-1}$. Sprzężenie nie zmienia rzędu,
zatem działanie jest dobrze określone.
Drugie twierdzenie Sylowa mówi, że jest ono przechodnie.
Stabilizatorem podgrupy $P$ jest z definicji jej normalizator
$N_G(P)$. Z twierdzenia o orbicie i stabilizatorze dostajemy
$n_p(G)=[G:N_G(P)]$.
Ponieważ $P\le N_G(P)\le G$, możemy zastosować
wzór Lagrange'a dwukrotnie:
\[
 [G:P]=[G:N_G(P)]\,[N_G(P):P].
\]
Lewa strona wynosi $m$, zatem $n_p(G)$ dzieli $m$.

Pozostaje wykazać kongruencję modulo $p$.
Niech $P$ działa na $\operatorname{Syl}_p(G)$
przez sprzężenie. Każda orbita tego działania ma
rozmiar będący potęgą $p$, ponieważ jej rozmiar dzieli
$|P|=p^a$. Wobec tego liczba wszystkich podgrup Sylowa
jest przystająca modulo $p$ do liczby orbit jednoelementowych,
czyli do liczby podgrup Sylowa ustalanych przez każde $u\in P$.

Pokażemy, że jedyną taką podgrupą jest $P$.
Oczywiście $P$ jest ustalana przez wszystkie elementy $P$,
gdyż $uPu^{-1}=P$ dla $u\in P$.
Niech teraz $Q\in\operatorname{Syl}_p(G)$ będzie ustalana
przez całe $P$. Oznacza to $uQu^{-1}=Q$ dla każdego $u\in P$,
czyli $P\subseteq N_G(Q)$.
Połóżmy $N=N_G(Q)$. Grupa $Q$ jest normalna w $N$,
a jej rząd to $p^a$.
Ponieważ $Q\le N\le G$, na mocy twierdzenia Lagrange'a
mamy $p^a=|Q|\mid|N|$ oraz $|N|\mid|G|=p^a m$.
Ponieważ $p\nmid m$, największą potęgą $p$ dzielącą
$|N|$ jest dokładnie $p^a$.
A zatem $P$ i $Q$ są obie podgrupami $p$-Sylowa \emph{grupy $N$}.
Stosując drugie twierdzenie Sylowa do $N$, otrzymujemy
$n\in N$ z $nPn^{-1}=Q$.
Normalność $Q$ w $N$ daje $n^{-1}Qn=Q$.
Sprzęgając poprzednią równość przez $n^{-1}$,
uzyskujemy $P=n^{-1}Qn=Q$.
Tak więc jedynym punktem stałym działania $P$
na $\operatorname{Syl}_p(G)$ jest $P$.
Pozostałe orbity mają rozmiary podzielne przez $p$,
więc $n_p(G)\equiv1\pmod p$.
\end{proof}

\begin{wniosek}[Kiedy podgrupa Sylowa jest normalna?]
\label{aa:wn:normalna-sylow}
Niech $G$ będzie grupą skończoną, $p$ liczbą pierwszą
dzielącą $|G|$, a $P\in\operatorname{Syl}_p(G)$.
Wtedy $P\trianglelefteq G$ dokładnie wtedy, gdy $n_p(G)=1$.
\end{wniosek}
\begin{proof}
Jeśli $P$ jest normalna, to każde sprzężenie $gPg^{-1}$
jest równe $P$. Ponieważ wszystkie podgrupy Sylowa są sprzężone,
nie istnieje inna podgrupa Sylowa, więc $n_p(G)=1$.
Odwrotnie, jeżeli istnieje tylko jedna podgrupa Sylowa,
to każde $gPg^{-1}$, będące podgrupą Sylowa, musi być jej równe.
To właśnie definicja normalności $P$.
\end{proof}

\section{Przykłady: jak rzeczywiście korzystać z trzech twierdzeń}
\begin{przyklad}[Podgrupy Sylowa w $S_3$]
\label{aa:prz:sylow-s3}
Mamy $|S_3|=6$.
Dla $p=3$ liczba $n_3$ dzieli $2$ i spełnia
$n_3\equiv1\pmod3$. Jedyna możliwość to $n_3=1$.
Rzeczywiście podgrupa $A_3=\langle(1\,2\,3)\rangle$
jest normalna i ma rząd $3$.
Dla $p=2$ mamy $n_2\mid3$ oraz $n_2\equiv1\pmod2$.
Możliwości są dwie: $1$ albo $3$.
Wypiszmy podgrupy rzędu $2$:
$\langle(1\,2)\rangle$, $\langle(1\,3)\rangle$
i $\langle(2\,3)\rangle$.
Są trzy i są wzajemnie sprzężone, więc $n_2=3$.
To pokazuje, że same warunki liczbowe Sylowa mogą pozostawiać
więcej niż jedną możliwość; trzeba wykorzystać strukturę grupy.
\end{przyklad}

\begin{przyklad}[Podgrupy Sylowa w $S_4$ i podziały na pary]
Niech $G=S_4$, $|G|=24=2^3\cdot3$.
Podgrupa $3$-Sylowa ma rząd $3$, a każda zawiera dwa
nieidentycznościowe $3$-cykle. Ponieważ w $S_4$ takich
cykli jest $\binom43\cdot2=8$, a dwie różne podgrupy
rzędu $3$ nie mają wspólnego nieidentycznościowego elementu,
otrzymujemy $n_3(S_4)=8/2=4$.

Dla $p=2$ trzecie twierdzenie Sylowa daje
$n_2\mid3$ i $n_2\equiv1\pmod2$, więc $n_2=1$ albo $3$.
Rozważmy trzy podziały zbioru $\{1,2,3,4\}$ na dwie
nieuporządkowane pary:
$\{\{1,2\},\{3,4\}\}$,
$\{\{1,3\},\{2,4\}\}$ oraz
$\{\{1,4\},\{2,3\}\}$.
Stabilizator podziału $\{\{1,2\},\{3,4\}\}$ to zbiór
permutacji, które albo zachowują obie pary z osobna,
albo je zamieniają. W pierwszym przypadku mamy $2\cdot2=4$
permutacje: w każdej parze niezależnie wybieramy, czy jej
elementy zamienić. W drugim przypadku obraz elementu $1$
można wybrać spośród $3,4$, a obraz $2$ jest wtedy
drugim z tych dwóch elementów; podobnie obrazy $3,4$
można ustalić na dwa sposoby w parze $\{1,2\}$.
Daje to $2\cdot2=4$ kolejne permutacje.
Obie rodziny są rozłączne, toteż stabilizator ma dokładnie
osiem elementów i jest podgrupą $2$-Sylowa.
Stabilizatory pozostałych dwóch podziałów również mają
osiem elementów, ponieważ podziały są wzajemnie
przenoszone przez permutacje $S_4$, a ich stabilizatory
są sprzężone.
Stabilizatory są parami różne. Istotnie transpozycja
$(1\,2)$ należy do stabilizatora pierwszego podziału,
lecz przestawia dwa pozostałe podziały i nie należy do ich
stabilizatorów. Analogicznie $(1\,3)$ należy do stabilizatora
drugiego podziału, lecz nie należy do stabilizatorów pierwszego
i trzeciego, a $(1\,4)$ należy wyłącznie do stabilizatora
trzeciego podziału. Mamy więc trzy różne podgrupy
$2$-Sylowa, czyli $n_2\ge3$.
Z ograniczenia Sylowa $n_2\mid3$ wynika $n_2=3$.
Ten przykład pokazuje, jak przejść od samych warunków
liczbowych do faktycznego zliczania podgrup.
\end{przyklad}

\begin{przyklad}[Podgrupy Sylowa w $A_4$]
\label{aa:prz:sylow-a4}
Rząd $A_4$ wynosi $12=2^2\cdot3$.
Jej elementy to identyczność, osiem $3$-cykli
i trzy iloczyny dwóch rozłącznych transpozycji.
Każda podgrupa rzędu $3$ zawiera dokładnie dwa
nieidentycznościowe elementy, będące wzajemnie odwrotnymi
$3$-cyklami. Dwie różne takie podgrupy nie mogą współdzielić
nieidentycznościowego elementu: element rzędu $3$
generuje jednoznacznie swoją trójelementową podgrupę.
Osiem $3$-cykli dzieli się więc na cztery pary,
czyli $n_3(A_4)=8/2=4$.
Jest to zgodne z warunkami $n_3\mid4$ i $n_3\equiv1\pmod3$.

Rozważmy zbiór
$V=\{e,a,b,c\}$, gdzie
$a=(1\,2)(3\,4)$, $b=(1\,3)(2\,4)$ oraz
$c=(1\,4)(2\,3)$.
Mamy $a^2=b^2=c^2=e$, ponieważ każdy z tych elementów
jest iloczynem rozłącznych transpozycji.
Przy konwencji składania od prawej do lewej bezpośrednie
obliczenie obrazów $1,2,3,4$ daje
$ab=ba=c$, $bc=cb=a$ oraz $ca=ac=b$.
Na przykład $ab$ wysyła $1\mapsto4$, $4\mapsto1$,
$2\mapsto3$, $3\mapsto2$, a więc $ab=c$.
Wszystkie iloczyny elementów $V$ znów należą do $V$.
Ponieważ $V$ zawiera identyczność i odwrotność każdego
ze swych elementów, jest podgrupą $A_4$.
Funkcja $(\mathbb Z/2\mathbb Z)^2\to V$ określona wzorem
$([i],[j])\mapsto a^i b^j$ jest homomorfizmem,
ponieważ $a$ i $b$ komutują. Jej obrazy czterech par
to kolejno $e,a,b,c$, więc jest bijekcją.
Dlatego $V\cong(\mathbb Z/2\mathbb Z)^2$.
Jest normalna w $A_4$, ponieważ sprzężenie permutacją
zachowuje długości cykli i rozłączność, zatem przestawia
trzy nieidentycznościowe elementy $V$.
Podgrupa $V$ ma rząd $4=2^2$, więc jest podgrupą $2$-Sylowa.
Z normalności i wniosku o normalności Sylowa wynika $n_2(A_4)=1$.
Nie można jednak wyciągnąć stąd wniosku, że $A_4$ jest
przemienna: choć $V$ jest przemienna i normalna,
w $A_4$ znajdują się także $3$-cykle, które nie komutują
ze wszystkimi elementami $V$. Dla przykładu
$(1\,2\,3)(1\,2)(3\,4)(1\,2\,3)^{-1}
=(1\,4)(2\,3)\ne (1\,2)(3\,4)$.
Obliczenia te wyjaśniają, dlaczego podgrupy Sylowa
nie muszą być cykliczne, a ich liczby nie muszą być jednakowe.
\end{przyklad}

\begin{twierdzenie}[Grupy rzędu $pq$ przy braku szczególnej kongruencji]
\label{aa:tw:pq-cykliczna}
Niech $p<q$ będą liczbami pierwszymi i niech $G$ będzie grupą
rzędu $pq$. Jeśli $p\nmid(q-1)$, to $G$ jest cykliczna.
\end{twierdzenie}
\begin{proof}
Dla podgrup $q$-Sylowa mamy $n_q\mid p$ oraz
$n_q\equiv1\pmod q$. Ponieważ $n_q\in\{1,p\}$
i $1<p<q$, liczba $p$ nie jest przystająca do $1$ modulo $q$.
Zatem $n_q=1$, czyli jedyna podgrupa $Q$ rzędu $q$ jest normalna.
Dla podgrup $p$-Sylowa mamy $n_p\mid q$,
więc $n_p\in\{1,q\}$, i $n_p\equiv1\pmod p$.
Założenie $p\nmid(q-1)$ wyklucza $n_p=q$.
Zatem $n_p=1$, czyli podgrupa $P$ rzędu $p$ także jest normalna.

Podgrupy $P$ i $Q$ są cykliczne, ponieważ mają rzędy pierwsze.
Ich przekrój jest trywialny: rząd $P\cap Q$ dzieli jednocześnie
$p$ i $q$, a te są różnymi liczbami pierwszymi.
Niech $x\in P$, $y\in Q$. Rozważmy komutator
$[x,y]=xyx^{-1}y^{-1}$.
Normalność $P$ daje $yx^{-1}y^{-1}\in P$,
więc $[x,y]=x(yx^{-1}y^{-1})\in P$.
Normalność $Q$ daje $xyx^{-1}\in Q$,
więc $[x,y]=(xyx^{-1})y^{-1}\in Q$.
Stąd $[x,y]\in P\cap Q=\{e\}$, a więc $xy=yx$.
Każdy element $P$ komutuje z każdym elementem $Q$.

Wybierzmy generatory $x$ grupy $P$ i $y$ grupy $Q$.
Ponieważ komutują i mają rzędy $p,q$, element $xy$ ma
rząd $pq$. Istotnie, jeśli $(xy)^k=e$, to
$x^k=y^{-k}\in P\cap Q$, więc $x^k=y^k=e$.
Zatem $p\mid k$ i $q\mid k$, a ponieważ $p,q$ są względnie
pierwsze, $pq\mid k$ (z tożsamości Bézouta).
Z drugiej strony $(xy)^{pq}=x^{pq}y^{pq}=e$.
Najmniejszym dodatnim takim $k$ jest więc $pq$.
Podgrupa $\langle xy\rangle$ ma $pq=|G|$ elementów,
toteż jest równa $G$. Grupa jest cykliczna.
\end{proof}

\begin{przyklad}[Dlaczego warunek $p\nmid(q-1)$ jest istotny?]
Niech $q=7$ i $p=3$; wówczas $3\mid(7-1)$.
Zbudujemy konkretną nieabelową grupę rzędu $21$.
Dla $a\in\mathbb Z/7\mathbb Z$ i $i\in\{0,1,2\}$
rozważmy permutację zbioru $\mathbb Z/7\mathbb Z$
zadaną przez
$f_{a,i}(x)=2^i x+a\pmod7$.
Każda jest odwracalna, bo $2^i$ jest niezerowe modulo $7$.
Funkcje $f_{a,i}$ są parami różne: $a=f_{a,i}(0)$,
a $2^i=f_{a,i}(1)-f_{a,i}(0)$, przy czym $1,2,4$
są różnymi klasami modulo $7$.
Jest ich zatem $7\cdot3=21$.
Składanie daje
$f_{a,i}\circ f_{b,j}=f_{a+2^ib,\,i+j\ (\mathrm{mod}\ 3)}$,
ponieważ $2^3\equiv1\pmod7$.
Identycznością jest $f_{0,0}$, a odwrotność też należy
do tej rodziny: można wziąć
$f_{-2^{-i}a,\,-i\ (\mathrm{mod}\ 3)}$,
przy czym $2^{-i}$ oznacza odwrotność $2^i$ modulo $7$.
Zatem te $21$ funkcji tworzy grupę względem składania.
Nie jest abelowa:
$f_{0,1}\circ f_{1,0}=f_{2,1}$,
natomiast $f_{1,0}\circ f_{0,1}=f_{1,1}$.
Z tego powodu nie jest cykliczna.

Podgrupa translacji $T=\{f_{a,0}:a\in\mathbb Z/7\mathbb Z\}$
ma rząd $7$ i jest normalna: dla $f_{b,j}$ sprzężenie
$f_{b,j}\circ f_{a,0}\circ f_{b,j}^{-1}$ jest translacją
o przesunięcie $2^j a$.
Podgrupa $\langle f_{0,1}\rangle$ ma rząd $3$,
lecz nie jest normalna: sprzężenie przez $f_{1,0}$
wysyła $f_{0,1}$ na $f_{-1,1}$, co nie należy
do $\langle f_{0,1}\rangle=\{f_{0,0},f_{0,1},f_{0,2}\}$.
Warunki Sylowa dają $n_7=1$ oraz $n_3\in\{1,7\}$;
z nienormalności podgrupy rzędu $3$ wynika $n_3=7$.
\end{przyklad}

\begin{twierdzenie}[Nie istnieje grupa prosta rzędu $30$]
\label{aa:tw:30-not-simple}
Niech $G$ będzie grupą rzędu $30$. Wtedy $G$ ma
właściwą nietrywialną podgrupę normalną.
Inaczej mówiąc, żadna grupa rzędu $30$ nie jest prosta.
\end{twierdzenie}
\begin{proof}
Przypomnijmy, że grupa \emph{prosta} to grupa $G\ne\{e\}$,
której jedynymi normalnymi podgrupami są $\{e\}$ i $G$.
Załóżmy przeciwnie, że $G$ jest prosta.
Ponieważ $|G|=2\cdot3\cdot5$,
trzecie twierdzenie Sylowa daje
$n_5\mid6$, $n_5\equiv1\pmod5$;
stąd $n_5\in\{1,6\}$.
Gdyby $n_5=1$, podgrupa rzędu $5$ byłaby normalna,
co przeczy prostocie. Zatem $n_5=6$.
Analogicznie $n_3\mid10$, $n_3\equiv1\pmod3$,
więc $n_3\in\{1,10\}$ i prostota wymusza $n_3=10$.

Każda podgrupa rzędu $5$ jest cykliczna i ma cztery
elementy różne od jedynki, wszystkie rzędu $5$.
Dwie różne podgrupy rzędu $5$ nie mają wspólnego
elementu nieidentycznościowego: taki element sam generowałby
obie, co byłoby sprzecznością.
Zatem w $G$ jest co najmniej $6\cdot4=24$
elementów rzędu $5$.
Podobnie dziesięć podgrup rzędu $3$ dostarcza
$10\cdot2=20$ różnych elementów rzędu $3$.
Żaden element nie może mieć równocześnie rzędu $3$ i $5$,
więc obie rodziny są rozłączne.
Dostalibyśmy co najmniej $24+20=44$ nieidentycznościowe
elementy w grupie rzędu $30$, co jest niemożliwe.
Zatem $G$ nie jest prosta.
\end{proof}

\section{Dodatkowe klasyczne zastosowanie: podgrupy małego indeksu}
\begin{twierdzenie}[Normalność podgrupy o indeksie najmniejszego dzielnika pierwszego]
\label{aa:tw:index-minprime}
Niech $G$ będzie grupą skończoną, a $p$ najmniejszą liczbą
pierwszą dzielącą $|G|$. Jeżeli $H\le G$ i $[G:H]=p$,
to $H$ jest normalna w $G$.
\end{twierdzenie}
\begin{proof}
Niech $G$ działa na $p$ warstwach lewostronnych $G/H$.
Działanie wyznacza homomorfizm
$\rho:G\to S_p$ i jego obraz $K=\rho(G)$.
Ponieważ działanie na warstwach jest przechodnie,
$K$ działa przechodnio na zbiorze $p$ elementów.
Z twierdzenia o orbicie i stabilizatorze, zastosowanego do $K$,
wynika, że $p\mid|K|$.
Z pierwszego twierdzenia o izomorfizmie
$|K|=|G|/|\ker\rho|$, więc $|K|$ dzieli $|G|$.
Ponadto $K\le S_p$, stąd $|K|\mid p!$.

Jeśli liczba pierwsza $q$ dzieli $|K|$, to dzieli także
$|G|$, a z minimalności $p$ mamy $q\ge p$.
Z drugiej strony $q\mid p!$ implikuje $q\le p$.
Dlatego jedynym dzielnikiem pierwszym $|K|$ może być $p$.
W liczbie $p!$ czynnik $p$ występuje tylko raz:
jedyną jego wielokrotnością wśród $1,2,\ldots,p$
jest liczba $p$.
Ponieważ $p\mid |K|$ i $|K|\mid p!$,
otrzymujemy $|K|=p$.

Zatem $|\ker\rho|=|G|/p=|H|$.
Jądro działania na warstwach zawiera się w stabilizatorze
warstwy $H$, którym jest $H$.
Mamy więc $\ker\rho\subseteq H$ i równość liczebności,
zatem $\ker\rho=H$.
Jako jądro homomorfizmu $H$ jest normalna w $G$.
\end{proof}

\begin{uwaga}[Jak korzystać z twierdzeń bez nieuprawnionych odwrotności?]
Wnioski Sylowa mówią o istnieniu podgrup określonego rzędu,
ich sprzężoności oraz ograniczeniach liczbowych ich liczby.
Nie twierdzą, że istnieje podgrupa dowolnego rzędu dzielącego $|G|$,
ani nie klasyfikują wszystkich grup danego rzędu.
Przy stosowaniu warunków $n_p\mid m$ i $n_p\equiv1\pmod p$
należy najpierw wyznaczyć wszystkie możliwości arytmetyczne,
a dopiero potem rozstrzygać je za pomocą dalszych informacji
o konkretnych elementach lub działaniach grupy.
\end{uwaga}
